% Start preamble
\documentclass[12pt,a4paper]{article}
\usepackage{geometry}
 \geometry{
 a4paper,
 total={170mm,257mm},
 left=20mm,
 top=20mm,
 }
\usepackage[utf8]{inputenc}
\usepackage[T1]{fontenc}
\usepackage[pdftex]{graphicx}
\graphicspath{{./}}
\usepackage{enumitem}
\usepackage{pdfpages}
\usepackage{hyperref}
\usepackage{tikz}
\usepackage{attachfile}
\usepackage{epstopdf}
\usepackage{array}
\usepackage{multirow}
\usepackage{multicol}
\usepackage{float}
%\usepackage[table]{xcolor,colorbl}
\setlength{\textwidth}{17cm}
\setlength{\oddsidemargin}{-0.5cm}
\setlength{\evensidemargin}{-0.5cm}
%\setlenght{\headsep}{0cm}
\setlength\parindent{0pt}
%\setlength{\extrarowheight}{3pt}
\usepackage{listings}
%\usepackage{xcolor}

\input{arduinoLanguage.tex}
%%%%%% Counting oppgaves %%%%%%
 \newcount\questnum \questnum=0
 \def\oppgave{
            \advance\questnum by 1
	    \ifthenelse{\questnum>0\AND \questnum<9}
	    {
                \vskip 1cm
		\textbf{Oppgave}\hskip 5pt\the\questnum \hfill \hfill(6p)
		\vskip 3pt
		\hrule
	\vskip 0.5cm}
	{
                \vskip 1cm
		\textbf{Oppgave}\hskip 5pt \the\questnum \hfill \hfill(12p)
		\vskip 3pt \hrule \vskip 0.5cm }

		}

\begin{document}
\title{TIF Prosesskjemi 01}
\author{Faglærer: Fred-Olav Mosdal 90507684\\}
\maketitle
\oppgave{}%1
\vskip 2.5pt 
På hvilke to måter koble klemmebrettet på en asynkronmotor? Lag en tegning som viser hvordan laskene legges i de to koblingene. 
\vskip 5pt 
\begin{tikzpicture}
	\draw[step=0.5cm,gray!90,very thin]  grid (17,8) ;
\end{tikzpicture}
\vskip 2.5pt 

\oppgave{}%2
\vskip 2.5pt 
Hva kaller vi den faste deler en i elektromotor, og hva kaller vi den roterende delen?
\vskip 5pt 
\begin{tikzpicture}
	\draw[step=0.5cm,gray!90,very thin]  grid (17,3) ;
\end{tikzpicture}
\vskip 2.5pt 
\oppgave{}%3
\vskip 2.5pt 
Hva mener vi med et «pådrag» i en prosess?
\vskip 5pt 
\begin{tikzpicture}
	\draw[step=0.5cm,gray!90,very thin]  grid (17,3);
\end{tikzpicture}
\vskip 2.5pt 
\newpage
\oppgave{}%4
\vskip 2.5pt 
Hva mener vi med en «utgang» i en prosess?
\vskip 5pt 
\begin{tikzpicture}
	\draw[step=0.5cm,gray!90,very thin]  grid (17,333) ;
\end{tikzpicture}
\vskip 2.5pt 
\oppgave{}%5
\vskip 2.5pt 
Hva er skal-verdi, er-verdi og avvik?
\vskip 5pt 
\begin{tikzpicture}
	\draw[step=0.5cm,gray!90,very thin]  grid (17,3) ;
\end{tikzpicture}
\vskip 2.5pt 
\oppgave{}%6
\vskip 2.5pt 
Hvordan er en reguleringssløyfe bygd opp?
\vskip 5pt 
\begin{tikzpicture}
	\draw[step=0.5cm,gray!90,very thin]  grid (17,3) ;
\end{tikzpicture}
\vskip 2.5pt 
\newpage
\oppgave{}%7
\vskip 2.5pt 
Du kjører i 45 km/h på en vei der fartsgrensen er 30 km/h. Hvor stort er avviket? Hva betyr det at du får et negativt svar?
\vskip 5pt 
\begin{tikzpicture}
	\draw[step=0.5cm,gray!90,very thin]  grid (17,6) ;
\end{tikzpicture}
\vskip 2.5pt 
\oppgave{}%8
\vskip 2.5pt 
Forklar hvordan et generelt måleinstrument er bygd opp.
\vskip 5pt 
\begin{tikzpicture}
	\draw[step=0.5cm,gray!90,very thin]  grid (17,3) ;
\end{tikzpicture}
\vskip 2.5pt 
\newpage
\oppgave{}%9
\vskip 2.5pt 
På figuren ser du en trykkmåler som er plassert i bunnen av en tank. Her er Pa trykket i toppen av tanken, p er tettheten til væsken, og h er væskehøyden målt fra bunnen av tanken. Hva er trykket i bunnen av tanken? Er det høyere, lavere eller det samme som i toppen av tanken?\\
\vskip 5pt 
$$\includegraphics[width=0.4\textwidth]{../output/nogpl/Trykkmåling.png}$$
\vskip 5pt 
\begin{tikzpicture}
	\draw[step=0.5cm,gray!90,very thin]  grid (17,19) ;
\end{tikzpicture}
\vskip 2.5pt 
\newpage
\end{document}
